\scanpage{22}
\mspage{9}

Donc pour prouver $2^\circ\Rightarrow3^\circ$, il \unsure{revient|reviendra} à prouver que $2^\circ$ \deleted{implique $4^\circ$} suffit à impliquer $4^\circ$ (et ceci quel que soit $I$ \inserted{\unsure{ordonné|considéré}} loc. fini) --- car en appliquant ceci aux $I_{<i_0}$, on trouve que $2^\circ$ implique $\alpha_{i_0}\in TC=C\cap W$ (non seulement $\alpha_{i_0}\in C$, selon la prop 2), donc $\beta_{i_0}\in TC\subset W$, donc $\widetilde u_{i_0}\in W$ par saturation, donc $\widetilde u_{i_0}\in TC$, d’où $3^\circ$.

Notons d’autre part que le b) de la prop.~2 \deleted{\unsure{n’est autre}} est conséquence \unsure{du|de} a), appliqué \deleted{\unsure{aux}} à la restriction de \unsure{d$^{\circ}$.|$u_\bullet$} à $I_{<i_0}$. Au total, il y a donc à prouver \unsure{seulement|seulem$^{\rm t}$} les deux choses

\[
\left\{\begin{aligned}
 \text{a)}\;&u_{\deleted{\unsure{I}}\bullet}\in C_I\Longrightarrow u_I\in C\\
 \text{b)}\;&u_\bullet\in C_I\cap W\Longrightarrow u_I\in C\cap W\quad(\overset{\text{déf}}{=}TC)\\
 &\hspace{5em}\inserted{[\text{et }u_\bullet\in(TC)_{\unsure{W|w}}]}
\end{aligned}\right. \tag*{(*)}
\]

Nous le prouverons par réduction \deleted{\unsure{à la}} aux cas particuliers de la prop.~1, déjà prouvée. J’ai d’ailleurs \unsure{aussi|ici} envie de prouver le

\uline{Corollaire 2} Les hypothèses \inserted{$\unsure{\mathcal M|\mathscr M},C$ (\deleted{\unsure{et}} \unsure{évt.|évent.} $W$),} sur $X_\bullet,Y_\bullet\in\mathcal M(I^\circ)_C$ et sur $u_\bullet\in C_I$, \inserted{\unsure{enfin|et} sur $I$ loc. fini,} sont celles de la prop.~2. \deleted{Soient} \inserted{Pour \unsure{tte} \inserted{partie} $J$ de $I$, soit} \deleted{$I'\subset I''$ deux \unsure{ouvertes|parties}}

\[
 \deleted{\unsure{X_J}}u_J\overset{\text{déf}}{=}\underset{i\in J}{\underrightarrow{\lim}}\,u_i:
 X_J\overset{\text{déf}}{=}\underset{\unsure{i'|i}\in J}{\underrightarrow{\lim}}\,X_i
 \mathbin{\unsure{-|\to}} Y_J\overset{\text{déf}}{=}\underset{i\in J}{\underrightarrow{\lim}}\,Y_i .
\]

Soient $I'\unsure{\subseteq|\subset} I''$ deux parties \uline{ouvertes} de $I$, et considérons \unsure{le|les} \unsure{diagramme|diagrammes}
% archive sheet number: 21

